% =====================================================================
% CHUONG 2 - CO SO LY THUYET
% Tac gia: Bui Dinh Manh (Evaluation Engineer & Documentation)
% Cach dung: them % =====================================================================
% CHUONG 2 - CO SO LY THUYET
% Tac gia: Bui Dinh Manh (Evaluation Engineer & Documentation)
% Cach dung: them % =====================================================================
% CHUONG 2 - CO SO LY THUYET
% Tac gia: Bui Dinh Manh (Evaluation Engineer & Documentation)
% Cach dung: them % =====================================================================
% CHUONG 2 - CO SO LY THUYET
% Tac gia: Bui Dinh Manh (Evaluation Engineer & Documentation)
% Cach dung: them \input{chuong2} vao main.tex.
% Neu dung documentclass{article}, doi \chapter -> \section.
% =====================================================================

\chapter{Cơ sở lý thuyết}
\label{chap:theory}

\section{Trí tuệ nhân tạo (Artificial Intelligence -- AI)}
\textbf{Trí tuệ nhân tạo} là lĩnh vực của khoa học máy tính nghiên cứu việc tạo
ra các hệ thống có khả năng thực hiện những nhiệm vụ vốn đòi hỏi trí thông minh
của con người, như nhận dạng hình ảnh, hiểu ngôn ngữ, ra quyết định và dự đoán.
AI bao gồm nhiều nhánh, trong đó \textbf{Học máy (Machine Learning)} là nhánh
phát triển mạnh mẽ và được ứng dụng rộng rãi nhất hiện nay.

\section{Học máy (Machine Learning -- ML)}
\textbf{Học máy} là phương pháp giúp máy tính "học" quy luật từ dữ liệu thay vì
được lập trình cứng bằng các quy tắc cố định. Có ba nhóm chính:
\begin{itemize}
  \item \textbf{Học có giám sát (Supervised Learning):} học từ dữ liệu đã gán
        nhãn để dự đoán nhãn cho dữ liệu mới (phân loại, hồi quy). \emph{Bài toán
        của đề tài thuộc nhóm này -- phân loại nhị phân.}
  \item \textbf{Học không giám sát (Unsupervised Learning):} tìm cấu trúc ẩn
        trong dữ liệu chưa gán nhãn (phân cụm, giảm chiều).
  \item \textbf{Học tăng cường (Reinforcement Learning):} học qua tương tác và
        phần thưởng.
\end{itemize}
Bài toán đánh giá nguy cơ bệnh tim là bài toán \textbf{phân loại nhị phân}
(binary classification): đầu ra là 1 (có nguy cơ) hoặc 0 (không có nguy cơ).

\section{Mạng nơ-ron nhân tạo (Artificial Neural Network -- ANN)}
\textbf{Mạng nơ-ron nhân tạo} là mô hình tính toán mô phỏng cách hoạt động của
các nơ-ron thần kinh trong não người. ANN gồm các \textbf{nơ-ron (neuron)} được
tổ chức thành nhiều \textbf{tầng (layer)}:
\begin{itemize}
  \item \textbf{Tầng đầu vào (Input Layer):} nhận vector đặc trưng (ở đề tài là
        13 đặc trưng).
  \item \textbf{Các tầng ẩn (Hidden Layers):} trích xuất và biến đổi đặc trưng;
        càng nhiều tầng/nơ-ron, mạng càng học được quan hệ phức tạp.
  \item \textbf{Tầng đầu ra (Output Layer):} đưa ra kết quả dự đoán (1 nơ-ron
        với hàm Sigmoid cho phân loại nhị phân).
\end{itemize}

Mỗi kết nối giữa hai nơ-ron có một \textbf{trọng số (weight)}, mỗi nơ-ron có
thêm một \textbf{độ lệch (bias)}. Giá trị tại một nơ-ron được tính bằng:
\begin{equation}
  z = \sum_{i=1}^{n} w_i x_i + b, \qquad a = f(z)
\end{equation}
trong đó $f$ là \textbf{hàm kích hoạt}. Quá trình tính toán từ đầu vào đến đầu ra
gọi là \textbf{lan truyền tiến (Forward Propagation)}.

\section{Hàm kích hoạt (Activation Function)}
Hàm kích hoạt đưa tính \textbf{phi tuyến} vào mạng, giúp ANN học được các quan
hệ phức tạp. Hai hàm dùng trong đề tài:
\begin{itemize}
  \item \textbf{ReLU (Rectified Linear Unit):} $f(z) = \max(0, z)$. Dùng ở các
        tầng ẩn vì tính toán nhanh, giảm hiện tượng \emph{vanishing gradient} so
        với sigmoid/tanh.
  \item \textbf{Sigmoid:} $f(z) = \dfrac{1}{1 + e^{-z}}$, cho giá trị trong
        khoảng $(0, 1)$. Dùng ở tầng đầu ra để biểu diễn \textbf{xác suất}
        thuộc lớp "có nguy cơ"; ngưỡng 0.5 dùng để quy về nhãn 0/1.
\end{itemize}

\section{Hàm mất mát (Loss Function)}
\textbf{Hàm mất mát} đo độ sai lệch giữa dự đoán và nhãn thực tế. Với phân loại
nhị phân, dùng \textbf{Binary Crossentropy (Log Loss)}:
\begin{equation}
  L = -\frac{1}{N} \sum_{i=1}^{N} \Big[ y_i \log(\hat{y}_i) + (1-y_i)\log(1-\hat{y}_i) \Big]
\end{equation}
trong đó $y$ là nhãn thực, $\hat{y}$ là xác suất dự đoán. Giá trị loss càng nhỏ
thì mô hình dự đoán càng chính xác.

\section{Lan truyền ngược (Backpropagation)}
\textbf{Backpropagation} là thuật toán cốt lõi để huấn luyện ANN. Ý tưởng:
\begin{enumerate}
  \item \textbf{Forward:} tính đầu ra và giá trị loss.
  \item \textbf{Backward:} tính \textbf{gradient} (đạo hàm) của loss theo từng
        trọng số, lan truyền ngược từ tầng đầu ra về tầng đầu vào theo quy tắc
        chuỗi (chain rule).
  \item \textbf{Cập nhật trọng số} theo hướng giảm loss (gradient descent):
        \begin{equation}
          w \leftarrow w - \eta \cdot \frac{\partial L}{\partial w}
        \end{equation}
        với $\eta$ là \textbf{tốc độ học (learning rate)}.
\end{enumerate}
Quá trình lặp lại qua nhiều \textbf{epoch} (lượt duyệt toàn bộ dữ liệu) cho đến
khi mô hình hội tụ.

\section{Bộ tối ưu (Optimizer) -- Adam}
\textbf{Adam (Adaptive Moment Estimation)} là bộ tối ưu cải tiến của gradient
descent, kết hợp \textbf{momentum} và \textbf{điều chỉnh tốc độ học thích nghi}
cho từng tham số. Adam hội tụ nhanh, ổn định và ít phải tinh chỉnh thủ công, nên
được chọn cho mô hình. Tốc độ học sử dụng là $0.001$, kích thước batch $32$.

\section{Quá khớp (Overfitting) và kỹ thuật chống quá khớp}
\textbf{Overfitting} xảy ra khi mô hình học "thuộc lòng" dữ liệu train mà kém
tổng quát trên dữ liệu mới (loss train thấp nhưng loss validation cao). Các kỹ
thuật chống:
\begin{itemize}
  \item \textbf{Early Stopping:} theo dõi loss trên tập validation, dừng huấn
        luyện khi loss không cải thiện sau một số epoch (patience) và giữ lại
        trọng số tốt nhất.
  \item \textbf{Regularization (L2):} thêm thành phần phạt vào loss để hạn chế
        trọng số lớn.
  \item \textbf{Dropout:} ngẫu nhiên "tắt" một phần nơ-ron khi huấn luyện
        (tuỳ chọn).
  \item \textbf{Chia tập dữ liệu hợp lý} và \textbf{loại bỏ dữ liệu trùng} để
        tránh rò rỉ dữ liệu.
\end{itemize}

\section{Các chỉ số đánh giá mô hình phân loại}
Dựa trên \textbf{ma trận nhầm lẫn (Confusion Matrix)} với 4 đại lượng: TP (dự
đoán đúng có bệnh), TN (đúng không bệnh), FP (báo nhầm có bệnh), FN (bỏ sót ca
có bệnh).
\begin{itemize}
  \item \textbf{Accuracy} $= \dfrac{TP + TN}{\text{Tổng}}$ -- tỷ lệ dự đoán đúng
        tổng thể.
  \item \textbf{Precision} $= \dfrac{TP}{TP + FP}$ -- độ chính xác của dự đoán
        "có bệnh".
  \item \textbf{Recall (Sensitivity)} $= \dfrac{TP}{TP + FN}$ -- khả năng phát
        hiện ca có bệnh (đặc biệt quan trọng trong y tế, vì bỏ sót bệnh nhân rất
        nguy hiểm).
  \item \textbf{F1-Score} $= \dfrac{2 \cdot Precision \cdot Recall}{Precision +
        Recall}$ -- trung bình điều hòa, cân bằng giữa Precision và Recall.
  \item \textbf{ROC-AUC:} diện tích dưới đường ROC, đo khả năng phân tách hai
        lớp ở mọi ngưỡng (0.5 = ngẫu nhiên, 1.0 = hoàn hảo).
\end{itemize}
Các chỉ số này được dùng để đánh giá và so sánh các mô hình ở \textbf{Chương 5}.
 vao main.tex.
% Neu dung documentclass{article}, doi \chapter -> \section.
% =====================================================================

\chapter{Cơ sở lý thuyết}
\label{chap:theory}

\section{Trí tuệ nhân tạo (Artificial Intelligence -- AI)}
\textbf{Trí tuệ nhân tạo} là lĩnh vực của khoa học máy tính nghiên cứu việc tạo
ra các hệ thống có khả năng thực hiện những nhiệm vụ vốn đòi hỏi trí thông minh
của con người, như nhận dạng hình ảnh, hiểu ngôn ngữ, ra quyết định và dự đoán.
AI bao gồm nhiều nhánh, trong đó \textbf{Học máy (Machine Learning)} là nhánh
phát triển mạnh mẽ và được ứng dụng rộng rãi nhất hiện nay.

\section{Học máy (Machine Learning -- ML)}
\textbf{Học máy} là phương pháp giúp máy tính "học" quy luật từ dữ liệu thay vì
được lập trình cứng bằng các quy tắc cố định. Có ba nhóm chính:
\begin{itemize}
  \item \textbf{Học có giám sát (Supervised Learning):} học từ dữ liệu đã gán
        nhãn để dự đoán nhãn cho dữ liệu mới (phân loại, hồi quy). \emph{Bài toán
        của đề tài thuộc nhóm này -- phân loại nhị phân.}
  \item \textbf{Học không giám sát (Unsupervised Learning):} tìm cấu trúc ẩn
        trong dữ liệu chưa gán nhãn (phân cụm, giảm chiều).
  \item \textbf{Học tăng cường (Reinforcement Learning):} học qua tương tác và
        phần thưởng.
\end{itemize}
Bài toán đánh giá nguy cơ bệnh tim là bài toán \textbf{phân loại nhị phân}
(binary classification): đầu ra là 1 (có nguy cơ) hoặc 0 (không có nguy cơ).

\section{Mạng nơ-ron nhân tạo (Artificial Neural Network -- ANN)}
\textbf{Mạng nơ-ron nhân tạo} là mô hình tính toán mô phỏng cách hoạt động của
các nơ-ron thần kinh trong não người. ANN gồm các \textbf{nơ-ron (neuron)} được
tổ chức thành nhiều \textbf{tầng (layer)}:
\begin{itemize}
  \item \textbf{Tầng đầu vào (Input Layer):} nhận vector đặc trưng (ở đề tài là
        13 đặc trưng).
  \item \textbf{Các tầng ẩn (Hidden Layers):} trích xuất và biến đổi đặc trưng;
        càng nhiều tầng/nơ-ron, mạng càng học được quan hệ phức tạp.
  \item \textbf{Tầng đầu ra (Output Layer):} đưa ra kết quả dự đoán (1 nơ-ron
        với hàm Sigmoid cho phân loại nhị phân).
\end{itemize}

Mỗi kết nối giữa hai nơ-ron có một \textbf{trọng số (weight)}, mỗi nơ-ron có
thêm một \textbf{độ lệch (bias)}. Giá trị tại một nơ-ron được tính bằng:
\begin{equation}
  z = \sum_{i=1}^{n} w_i x_i + b, \qquad a = f(z)
\end{equation}
trong đó $f$ là \textbf{hàm kích hoạt}. Quá trình tính toán từ đầu vào đến đầu ra
gọi là \textbf{lan truyền tiến (Forward Propagation)}.

\section{Hàm kích hoạt (Activation Function)}
Hàm kích hoạt đưa tính \textbf{phi tuyến} vào mạng, giúp ANN học được các quan
hệ phức tạp. Hai hàm dùng trong đề tài:
\begin{itemize}
  \item \textbf{ReLU (Rectified Linear Unit):} $f(z) = \max(0, z)$. Dùng ở các
        tầng ẩn vì tính toán nhanh, giảm hiện tượng \emph{vanishing gradient} so
        với sigmoid/tanh.
  \item \textbf{Sigmoid:} $f(z) = \dfrac{1}{1 + e^{-z}}$, cho giá trị trong
        khoảng $(0, 1)$. Dùng ở tầng đầu ra để biểu diễn \textbf{xác suất}
        thuộc lớp "có nguy cơ"; ngưỡng 0.5 dùng để quy về nhãn 0/1.
\end{itemize}

\section{Hàm mất mát (Loss Function)}
\textbf{Hàm mất mát} đo độ sai lệch giữa dự đoán và nhãn thực tế. Với phân loại
nhị phân, dùng \textbf{Binary Crossentropy (Log Loss)}:
\begin{equation}
  L = -\frac{1}{N} \sum_{i=1}^{N} \Big[ y_i \log(\hat{y}_i) + (1-y_i)\log(1-\hat{y}_i) \Big]
\end{equation}
trong đó $y$ là nhãn thực, $\hat{y}$ là xác suất dự đoán. Giá trị loss càng nhỏ
thì mô hình dự đoán càng chính xác.

\section{Lan truyền ngược (Backpropagation)}
\textbf{Backpropagation} là thuật toán cốt lõi để huấn luyện ANN. Ý tưởng:
\begin{enumerate}
  \item \textbf{Forward:} tính đầu ra và giá trị loss.
  \item \textbf{Backward:} tính \textbf{gradient} (đạo hàm) của loss theo từng
        trọng số, lan truyền ngược từ tầng đầu ra về tầng đầu vào theo quy tắc
        chuỗi (chain rule).
  \item \textbf{Cập nhật trọng số} theo hướng giảm loss (gradient descent):
        \begin{equation}
          w \leftarrow w - \eta \cdot \frac{\partial L}{\partial w}
        \end{equation}
        với $\eta$ là \textbf{tốc độ học (learning rate)}.
\end{enumerate}
Quá trình lặp lại qua nhiều \textbf{epoch} (lượt duyệt toàn bộ dữ liệu) cho đến
khi mô hình hội tụ.

\section{Bộ tối ưu (Optimizer) -- Adam}
\textbf{Adam (Adaptive Moment Estimation)} là bộ tối ưu cải tiến của gradient
descent, kết hợp \textbf{momentum} và \textbf{điều chỉnh tốc độ học thích nghi}
cho từng tham số. Adam hội tụ nhanh, ổn định và ít phải tinh chỉnh thủ công, nên
được chọn cho mô hình. Tốc độ học sử dụng là $0.001$, kích thước batch $32$.

\section{Quá khớp (Overfitting) và kỹ thuật chống quá khớp}
\textbf{Overfitting} xảy ra khi mô hình học "thuộc lòng" dữ liệu train mà kém
tổng quát trên dữ liệu mới (loss train thấp nhưng loss validation cao). Các kỹ
thuật chống:
\begin{itemize}
  \item \textbf{Early Stopping:} theo dõi loss trên tập validation, dừng huấn
        luyện khi loss không cải thiện sau một số epoch (patience) và giữ lại
        trọng số tốt nhất.
  \item \textbf{Regularization (L2):} thêm thành phần phạt vào loss để hạn chế
        trọng số lớn.
  \item \textbf{Dropout:} ngẫu nhiên "tắt" một phần nơ-ron khi huấn luyện
        (tuỳ chọn).
  \item \textbf{Chia tập dữ liệu hợp lý} và \textbf{loại bỏ dữ liệu trùng} để
        tránh rò rỉ dữ liệu.
\end{itemize}

\section{Các chỉ số đánh giá mô hình phân loại}
Dựa trên \textbf{ma trận nhầm lẫn (Confusion Matrix)} với 4 đại lượng: TP (dự
đoán đúng có bệnh), TN (đúng không bệnh), FP (báo nhầm có bệnh), FN (bỏ sót ca
có bệnh).
\begin{itemize}
  \item \textbf{Accuracy} $= \dfrac{TP + TN}{\text{Tổng}}$ -- tỷ lệ dự đoán đúng
        tổng thể.
  \item \textbf{Precision} $= \dfrac{TP}{TP + FP}$ -- độ chính xác của dự đoán
        "có bệnh".
  \item \textbf{Recall (Sensitivity)} $= \dfrac{TP}{TP + FN}$ -- khả năng phát
        hiện ca có bệnh (đặc biệt quan trọng trong y tế, vì bỏ sót bệnh nhân rất
        nguy hiểm).
  \item \textbf{F1-Score} $= \dfrac{2 \cdot Precision \cdot Recall}{Precision +
        Recall}$ -- trung bình điều hòa, cân bằng giữa Precision và Recall.
  \item \textbf{ROC-AUC:} diện tích dưới đường ROC, đo khả năng phân tách hai
        lớp ở mọi ngưỡng (0.5 = ngẫu nhiên, 1.0 = hoàn hảo).
\end{itemize}
Các chỉ số này được dùng để đánh giá và so sánh các mô hình ở \textbf{Chương 5}.
 vao main.tex.
% Neu dung documentclass{article}, doi \chapter -> \section.
% =====================================================================

\chapter{Cơ sở lý thuyết}
\label{chap:theory}

\section{Trí tuệ nhân tạo (Artificial Intelligence -- AI)}
\textbf{Trí tuệ nhân tạo} là lĩnh vực của khoa học máy tính nghiên cứu việc tạo
ra các hệ thống có khả năng thực hiện những nhiệm vụ vốn đòi hỏi trí thông minh
của con người, như nhận dạng hình ảnh, hiểu ngôn ngữ, ra quyết định và dự đoán.
AI bao gồm nhiều nhánh, trong đó \textbf{Học máy (Machine Learning)} là nhánh
phát triển mạnh mẽ và được ứng dụng rộng rãi nhất hiện nay.

\section{Học máy (Machine Learning -- ML)}
\textbf{Học máy} là phương pháp giúp máy tính "học" quy luật từ dữ liệu thay vì
được lập trình cứng bằng các quy tắc cố định. Có ba nhóm chính:
\begin{itemize}
  \item \textbf{Học có giám sát (Supervised Learning):} học từ dữ liệu đã gán
        nhãn để dự đoán nhãn cho dữ liệu mới (phân loại, hồi quy). \emph{Bài toán
        của đề tài thuộc nhóm này -- phân loại nhị phân.}
  \item \textbf{Học không giám sát (Unsupervised Learning):} tìm cấu trúc ẩn
        trong dữ liệu chưa gán nhãn (phân cụm, giảm chiều).
  \item \textbf{Học tăng cường (Reinforcement Learning):} học qua tương tác và
        phần thưởng.
\end{itemize}
Bài toán đánh giá nguy cơ bệnh tim là bài toán \textbf{phân loại nhị phân}
(binary classification): đầu ra là 1 (có nguy cơ) hoặc 0 (không có nguy cơ).

\section{Mạng nơ-ron nhân tạo (Artificial Neural Network -- ANN)}
\textbf{Mạng nơ-ron nhân tạo} là mô hình tính toán mô phỏng cách hoạt động của
các nơ-ron thần kinh trong não người. ANN gồm các \textbf{nơ-ron (neuron)} được
tổ chức thành nhiều \textbf{tầng (layer)}:
\begin{itemize}
  \item \textbf{Tầng đầu vào (Input Layer):} nhận vector đặc trưng (ở đề tài là
        13 đặc trưng).
  \item \textbf{Các tầng ẩn (Hidden Layers):} trích xuất và biến đổi đặc trưng;
        càng nhiều tầng/nơ-ron, mạng càng học được quan hệ phức tạp.
  \item \textbf{Tầng đầu ra (Output Layer):} đưa ra kết quả dự đoán (1 nơ-ron
        với hàm Sigmoid cho phân loại nhị phân).
\end{itemize}

Mỗi kết nối giữa hai nơ-ron có một \textbf{trọng số (weight)}, mỗi nơ-ron có
thêm một \textbf{độ lệch (bias)}. Giá trị tại một nơ-ron được tính bằng:
\begin{equation}
  z = \sum_{i=1}^{n} w_i x_i + b, \qquad a = f(z)
\end{equation}
trong đó $f$ là \textbf{hàm kích hoạt}. Quá trình tính toán từ đầu vào đến đầu ra
gọi là \textbf{lan truyền tiến (Forward Propagation)}.

\section{Hàm kích hoạt (Activation Function)}
Hàm kích hoạt đưa tính \textbf{phi tuyến} vào mạng, giúp ANN học được các quan
hệ phức tạp. Hai hàm dùng trong đề tài:
\begin{itemize}
  \item \textbf{ReLU (Rectified Linear Unit):} $f(z) = \max(0, z)$. Dùng ở các
        tầng ẩn vì tính toán nhanh, giảm hiện tượng \emph{vanishing gradient} so
        với sigmoid/tanh.
  \item \textbf{Sigmoid:} $f(z) = \dfrac{1}{1 + e^{-z}}$, cho giá trị trong
        khoảng $(0, 1)$. Dùng ở tầng đầu ra để biểu diễn \textbf{xác suất}
        thuộc lớp "có nguy cơ"; ngưỡng 0.5 dùng để quy về nhãn 0/1.
\end{itemize}

\section{Hàm mất mát (Loss Function)}
\textbf{Hàm mất mát} đo độ sai lệch giữa dự đoán và nhãn thực tế. Với phân loại
nhị phân, dùng \textbf{Binary Crossentropy (Log Loss)}:
\begin{equation}
  L = -\frac{1}{N} \sum_{i=1}^{N} \Big[ y_i \log(\hat{y}_i) + (1-y_i)\log(1-\hat{y}_i) \Big]
\end{equation}
trong đó $y$ là nhãn thực, $\hat{y}$ là xác suất dự đoán. Giá trị loss càng nhỏ
thì mô hình dự đoán càng chính xác.

\section{Lan truyền ngược (Backpropagation)}
\textbf{Backpropagation} là thuật toán cốt lõi để huấn luyện ANN. Ý tưởng:
\begin{enumerate}
  \item \textbf{Forward:} tính đầu ra và giá trị loss.
  \item \textbf{Backward:} tính \textbf{gradient} (đạo hàm) của loss theo từng
        trọng số, lan truyền ngược từ tầng đầu ra về tầng đầu vào theo quy tắc
        chuỗi (chain rule).
  \item \textbf{Cập nhật trọng số} theo hướng giảm loss (gradient descent):
        \begin{equation}
          w \leftarrow w - \eta \cdot \frac{\partial L}{\partial w}
        \end{equation}
        với $\eta$ là \textbf{tốc độ học (learning rate)}.
\end{enumerate}
Quá trình lặp lại qua nhiều \textbf{epoch} (lượt duyệt toàn bộ dữ liệu) cho đến
khi mô hình hội tụ.

\section{Bộ tối ưu (Optimizer) -- Adam}
\textbf{Adam (Adaptive Moment Estimation)} là bộ tối ưu cải tiến của gradient
descent, kết hợp \textbf{momentum} và \textbf{điều chỉnh tốc độ học thích nghi}
cho từng tham số. Adam hội tụ nhanh, ổn định và ít phải tinh chỉnh thủ công, nên
được chọn cho mô hình. Tốc độ học sử dụng là $0.001$, kích thước batch $32$.

\section{Quá khớp (Overfitting) và kỹ thuật chống quá khớp}
\textbf{Overfitting} xảy ra khi mô hình học "thuộc lòng" dữ liệu train mà kém
tổng quát trên dữ liệu mới (loss train thấp nhưng loss validation cao). Các kỹ
thuật chống:
\begin{itemize}
  \item \textbf{Early Stopping:} theo dõi loss trên tập validation, dừng huấn
        luyện khi loss không cải thiện sau một số epoch (patience) và giữ lại
        trọng số tốt nhất.
  \item \textbf{Regularization (L2):} thêm thành phần phạt vào loss để hạn chế
        trọng số lớn.
  \item \textbf{Dropout:} ngẫu nhiên "tắt" một phần nơ-ron khi huấn luyện
        (tuỳ chọn).
  \item \textbf{Chia tập dữ liệu hợp lý} và \textbf{loại bỏ dữ liệu trùng} để
        tránh rò rỉ dữ liệu.
\end{itemize}

\section{Các chỉ số đánh giá mô hình phân loại}
Dựa trên \textbf{ma trận nhầm lẫn (Confusion Matrix)} với 4 đại lượng: TP (dự
đoán đúng có bệnh), TN (đúng không bệnh), FP (báo nhầm có bệnh), FN (bỏ sót ca
có bệnh).
\begin{itemize}
  \item \textbf{Accuracy} $= \dfrac{TP + TN}{\text{Tổng}}$ -- tỷ lệ dự đoán đúng
        tổng thể.
  \item \textbf{Precision} $= \dfrac{TP}{TP + FP}$ -- độ chính xác của dự đoán
        "có bệnh".
  \item \textbf{Recall (Sensitivity)} $= \dfrac{TP}{TP + FN}$ -- khả năng phát
        hiện ca có bệnh (đặc biệt quan trọng trong y tế, vì bỏ sót bệnh nhân rất
        nguy hiểm).
  \item \textbf{F1-Score} $= \dfrac{2 \cdot Precision \cdot Recall}{Precision +
        Recall}$ -- trung bình điều hòa, cân bằng giữa Precision và Recall.
  \item \textbf{ROC-AUC:} diện tích dưới đường ROC, đo khả năng phân tách hai
        lớp ở mọi ngưỡng (0.5 = ngẫu nhiên, 1.0 = hoàn hảo).
\end{itemize}
Các chỉ số này được dùng để đánh giá và so sánh các mô hình ở \textbf{Chương 5}.
 vao main.tex.
% Neu dung documentclass{article}, doi \chapter -> \section.
% =====================================================================

\chapter{Cơ sở lý thuyết}
\label{chap:theory}

\section{Trí tuệ nhân tạo (Artificial Intelligence -- AI)}
\textbf{Trí tuệ nhân tạo} là lĩnh vực của khoa học máy tính nghiên cứu việc tạo
ra các hệ thống có khả năng thực hiện những nhiệm vụ vốn đòi hỏi trí thông minh
của con người, như nhận dạng hình ảnh, hiểu ngôn ngữ, ra quyết định và dự đoán.
AI bao gồm nhiều nhánh, trong đó \textbf{Học máy (Machine Learning)} là nhánh
phát triển mạnh mẽ và được ứng dụng rộng rãi nhất hiện nay.

\section{Học máy (Machine Learning -- ML)}
\textbf{Học máy} là phương pháp giúp máy tính "học" quy luật từ dữ liệu thay vì
được lập trình cứng bằng các quy tắc cố định. Có ba nhóm chính:
\begin{itemize}
  \item \textbf{Học có giám sát (Supervised Learning):} học từ dữ liệu đã gán
        nhãn để dự đoán nhãn cho dữ liệu mới (phân loại, hồi quy). \emph{Bài toán
        của đề tài thuộc nhóm này -- phân loại nhị phân.}
  \item \textbf{Học không giám sát (Unsupervised Learning):} tìm cấu trúc ẩn
        trong dữ liệu chưa gán nhãn (phân cụm, giảm chiều).
  \item \textbf{Học tăng cường (Reinforcement Learning):} học qua tương tác và
        phần thưởng.
\end{itemize}
Bài toán đánh giá nguy cơ bệnh tim là bài toán \textbf{phân loại nhị phân}
(binary classification): đầu ra là 1 (có nguy cơ) hoặc 0 (không có nguy cơ).

\section{Mạng nơ-ron nhân tạo (Artificial Neural Network -- ANN)}
\textbf{Mạng nơ-ron nhân tạo} là mô hình tính toán mô phỏng cách hoạt động của
các nơ-ron thần kinh trong não người. ANN gồm các \textbf{nơ-ron (neuron)} được
tổ chức thành nhiều \textbf{tầng (layer)}:
\begin{itemize}
  \item \textbf{Tầng đầu vào (Input Layer):} nhận vector đặc trưng (ở đề tài là
        13 đặc trưng).
  \item \textbf{Các tầng ẩn (Hidden Layers):} trích xuất và biến đổi đặc trưng;
        càng nhiều tầng/nơ-ron, mạng càng học được quan hệ phức tạp.
  \item \textbf{Tầng đầu ra (Output Layer):} đưa ra kết quả dự đoán (1 nơ-ron
        với hàm Sigmoid cho phân loại nhị phân).
\end{itemize}

Mỗi kết nối giữa hai nơ-ron có một \textbf{trọng số (weight)}, mỗi nơ-ron có
thêm một \textbf{độ lệch (bias)}. Giá trị tại một nơ-ron được tính bằng:
\begin{equation}
  z = \sum_{i=1}^{n} w_i x_i + b, \qquad a = f(z)
\end{equation}
trong đó $f$ là \textbf{hàm kích hoạt}. Quá trình tính toán từ đầu vào đến đầu ra
gọi là \textbf{lan truyền tiến (Forward Propagation)}.

\section{Hàm kích hoạt (Activation Function)}
Hàm kích hoạt đưa tính \textbf{phi tuyến} vào mạng, giúp ANN học được các quan
hệ phức tạp. Hai hàm dùng trong đề tài:
\begin{itemize}
  \item \textbf{ReLU (Rectified Linear Unit):} $f(z) = \max(0, z)$. Dùng ở các
        tầng ẩn vì tính toán nhanh, giảm hiện tượng \emph{vanishing gradient} so
        với sigmoid/tanh.
  \item \textbf{Sigmoid:} $f(z) = \dfrac{1}{1 + e^{-z}}$, cho giá trị trong
        khoảng $(0, 1)$. Dùng ở tầng đầu ra để biểu diễn \textbf{xác suất}
        thuộc lớp "có nguy cơ"; ngưỡng 0.5 dùng để quy về nhãn 0/1.
\end{itemize}

\section{Hàm mất mát (Loss Function)}
\textbf{Hàm mất mát} đo độ sai lệch giữa dự đoán và nhãn thực tế. Với phân loại
nhị phân, dùng \textbf{Binary Crossentropy (Log Loss)}:
\begin{equation}
  L = -\frac{1}{N} \sum_{i=1}^{N} \Big[ y_i \log(\hat{y}_i) + (1-y_i)\log(1-\hat{y}_i) \Big]
\end{equation}
trong đó $y$ là nhãn thực, $\hat{y}$ là xác suất dự đoán. Giá trị loss càng nhỏ
thì mô hình dự đoán càng chính xác.

\section{Lan truyền ngược (Backpropagation)}
\textbf{Backpropagation} là thuật toán cốt lõi để huấn luyện ANN. Ý tưởng:
\begin{enumerate}
  \item \textbf{Forward:} tính đầu ra và giá trị loss.
  \item \textbf{Backward:} tính \textbf{gradient} (đạo hàm) của loss theo từng
        trọng số, lan truyền ngược từ tầng đầu ra về tầng đầu vào theo quy tắc
        chuỗi (chain rule).
  \item \textbf{Cập nhật trọng số} theo hướng giảm loss (gradient descent):
        \begin{equation}
          w \leftarrow w - \eta \cdot \frac{\partial L}{\partial w}
        \end{equation}
        với $\eta$ là \textbf{tốc độ học (learning rate)}.
\end{enumerate}
Quá trình lặp lại qua nhiều \textbf{epoch} (lượt duyệt toàn bộ dữ liệu) cho đến
khi mô hình hội tụ.

\section{Bộ tối ưu (Optimizer) -- Adam}
\textbf{Adam (Adaptive Moment Estimation)} là bộ tối ưu cải tiến của gradient
descent, kết hợp \textbf{momentum} và \textbf{điều chỉnh tốc độ học thích nghi}
cho từng tham số. Adam hội tụ nhanh, ổn định và ít phải tinh chỉnh thủ công, nên
được chọn cho mô hình. Tốc độ học sử dụng là $0.001$, kích thước batch $32$.

\section{Quá khớp (Overfitting) và kỹ thuật chống quá khớp}
\textbf{Overfitting} xảy ra khi mô hình học "thuộc lòng" dữ liệu train mà kém
tổng quát trên dữ liệu mới (loss train thấp nhưng loss validation cao). Các kỹ
thuật chống:
\begin{itemize}
  \item \textbf{Early Stopping:} theo dõi loss trên tập validation, dừng huấn
        luyện khi loss không cải thiện sau một số epoch (patience) và giữ lại
        trọng số tốt nhất.
  \item \textbf{Regularization (L2):} thêm thành phần phạt vào loss để hạn chế
        trọng số lớn.
  \item \textbf{Dropout:} ngẫu nhiên "tắt" một phần nơ-ron khi huấn luyện
        (tuỳ chọn).
  \item \textbf{Chia tập dữ liệu hợp lý} và \textbf{loại bỏ dữ liệu trùng} để
        tránh rò rỉ dữ liệu.
\end{itemize}

\section{Các chỉ số đánh giá mô hình phân loại}
Dựa trên \textbf{ma trận nhầm lẫn (Confusion Matrix)} với 4 đại lượng: TP (dự
đoán đúng có bệnh), TN (đúng không bệnh), FP (báo nhầm có bệnh), FN (bỏ sót ca
có bệnh).
\begin{itemize}
  \item \textbf{Accuracy} $= \dfrac{TP + TN}{\text{Tổng}}$ -- tỷ lệ dự đoán đúng
        tổng thể.
  \item \textbf{Precision} $= \dfrac{TP}{TP + FP}$ -- độ chính xác của dự đoán
        "có bệnh".
  \item \textbf{Recall (Sensitivity)} $= \dfrac{TP}{TP + FN}$ -- khả năng phát
        hiện ca có bệnh (đặc biệt quan trọng trong y tế, vì bỏ sót bệnh nhân rất
        nguy hiểm).
  \item \textbf{F1-Score} $= \dfrac{2 \cdot Precision \cdot Recall}{Precision +
        Recall}$ -- trung bình điều hòa, cân bằng giữa Precision và Recall.
  \item \textbf{ROC-AUC:} diện tích dưới đường ROC, đo khả năng phân tách hai
        lớp ở mọi ngưỡng (0.5 = ngẫu nhiên, 1.0 = hoàn hảo).
\end{itemize}
Các chỉ số này được dùng để đánh giá và so sánh các mô hình ở \textbf{Chương 5}.
